problem:  
Let \( M^n \) be a compact, simply connected Riemannian manifold of nonnegative sectional curvature admitting an **effective, isometric \( T^k \)-action** satisfying  
\[
k + \max_{x \in M} \dim(T_x) = n - 1,
\]
so that the action is *almost isotropy-maximal* (but not isotropy-maximal).  

**Problem (Conjectural Classification Question):**  
Is every such \( M \) diffeomorphic to a manifold that arises as a smooth quotient of an **isotropy-maximal** nonnegatively curved manifold by a freely acting circle subgroup of its isometry group?  
Equivalently, does relaxing isotropy-maximality by one correspond precisely to taking circle quotients (or small free extensions) of known isotropy‑maximal examples, such as products of spheres, compact rank-one symmetric spaces (CROSSes), or homogeneous biquotients of compact Lie groups?  

If not, can one construct genuinely new diffeomorphism types of nonnegatively curved manifolds that admit almost isotropy-maximal torus actions?  

---

Why is it a "good" problem:  

This problem is **precise, natural, and deeply tied to a known boundary in the field**. It asks whether the step from isotropy‑maximal to almost isotropy‑maximal actions corresponds to a simple topological operation (quotient by a circle action) or introduces fundamentally new examples—an unresolved question at the core of large‑symmetry nonnegative curvature theory.  

1. **Natural refinement of a known classification.** Classification results for isotropy‑maximal nonnegatively curved manifolds (by Grove–Searle, Wilking, Galaz‑García, Kerin, Searle, etc.) are complete. Asking what happens when we decrease isotropy rank by exactly one probes the immediate, unexplored frontier of these results.  

2. **Clear and narrowly defined conjectural focus.** The problem makes a concrete conjecture: *every almost isotropy‑maximal example arises from isotropy‑maximal ones by a circle quotient.* This can be proved or disproved, giving a definite yes/no statement suitable for research.  

3. **Deep cross-field connections.** Resolution would connect transformation group theory (torus actions and isotropy types), Riemannian geometry (nonnegative curvature, orbit metrics), and topology (circle quotients, biquotients, diffeomorphism classification).  

4. **Potential for conceptual breakthroughs.** Proving the conjecture would demonstrate that relaxing isotropy by one does not yield new topology—a rigidity phenomenon; disproving it would exhibit the first new topological types beyond known classification boundaries. Either result would have strong implications for understanding symmetry–curvature interplay.  

Thus, the problem is simple to state, mathematically precise, logically consistent, and strategically positioned where major progress could reveal new geometric principles—a hallmark of a truly “good” research problem.